\documentclass[11pt]{article}
\usepackage[a4paper,margin=1in]{geometry}
\usepackage{fontspec}
\usepackage[boldfont=false]{xeCJK}
\setCJKmainfont{FandolSong-Regular.otf}
\usepackage{amsmath}
\usepackage{amssymb}
\usepackage{array}
\usepackage{booktabs}
\usepackage{graphicx}
\usepackage{adjustbox}
\usepackage{enumitem}
\setlist[itemize]{topsep=2pt,partopsep=0pt,parsep=0pt,itemsep=2pt,leftmargin=1.4em}
\setlist[enumerate]{topsep=2pt,partopsep=0pt,parsep=0pt,itemsep=2pt,leftmargin=1.6em}
\usepackage{parskip}
\usepackage{fvextra}
\fvset{breaklines=true,breakanywhere=true,fontsize=\small}
\setlength{\parindent}{0pt}

\begin{document}

\begin{center}
  {\LARGE\bfseries Organisation of the organism}\\[0.35em]
  {\large Igcse Biology}
\end{center}
\vspace{0.6em}

\section*{Cells: the building blocks of life}

Every living thing is made of \textbf{cells} 细胞. A cell is the smallest part that can carry out the life processes. In this topic you compare a \textbf{plant} 植物 cell, an \textbf{animal} 动物 cell and the cell of a \textbf{bacterium} 细菌. All new cells are made by the \textbf{division} 分裂 of cells that already exist.

\subsection*{Which structures are in each cell?}

\begin{center}
\adjustbox{max width=\linewidth}{%
\begin{tabular}{llll}
\toprule
\textbf{Structure} & \textbf{Animal cell} & \textbf{Plant cell} & \textbf{Bacterial cell} \\
\midrule
\textbf{cell membrane} 细胞膜 & yes & yes & yes \\
\textbf{cytoplasm} 细胞质 & yes & yes & yes \\
\textbf{ribosomes} 核糖体 & yes & yes & yes \\
\textbf{nucleus} 细胞核 & yes & yes & no — has a loop of DNA instead \\
\textbf{mitochondria} 线粒体 & yes & yes & no \\
\textbf{cell wall} 细胞壁 & no & yes & yes (different material) \\
\textbf{chloroplasts} 叶绿体 & no & yes (only in green parts) & no \\
large \textbf{vacuole} 液泡 & no & yes & no \\
plasmids & no & no & yes \\
\bottomrule
\end{tabular}}
\end{center}

So plant and animal cells share five structures: cell membrane, cytoplasm, ribosomes, nucleus and mitochondria. A plant cell has three extra structures: a cell wall, a large vacuole and (in green parts) chloroplasts.

\subsection*{What each structure does}

\begin{itemize}
\item \textbf{Cell membrane} — controls which substances enter and leave the cell.

\item \textbf{Cytoplasm} — a jelly-like liquid where many chemical reactions happen.

\item \textbf{Nucleus} — controls the cell's activities and holds the \textbf{genetic material} 遗传物质.

\item \textbf{Ribosomes} — where \textbf{proteins} 蛋白质 are made.

\item \textbf{Mitochondria} — where \textbf{respiration} 呼吸作用 happens to release \textbf{energy} 能量.

\item \textbf{Cell wall} — made of \textbf{cellulose} 纤维素; it gives a plant cell strength and a fixed shape.

\item \textbf{Chloroplasts} — contain \textbf{chlorophyll} 叶绿素, a green substance that traps light for \textbf{photosynthesis} 光合作用.

\item \textbf{Large vacuole} — filled with \textbf{cell sap} 细胞液; it helps keep the plant cell firm.

\end{itemize}

\subsection*{Bacterial cells}

A \textbf{bacterium} is a \textbf{prokaryote} 原核生物 — its cell has \textbf{no} nucleus. Instead its DNA is a single loop (a ring), loose in the cytoplasm. Bacteria also have small extra rings of DNA called \textbf{plasmids} 质粒. A bacterial cell has a cell wall and cell membrane, cytoplasm and ribosomes, but no mitochondria and no chloroplasts.

\section*{Specialised cells}

Most cells are \textbf{specialised cells} 特化细胞 — their shape and parts suit one special job.

\begin{center}
\adjustbox{max width=\linewidth}{%
\begin{tabular}{lll}
\toprule
\textbf{Cell} & \textbf{Job} & \textbf{How its shape helps} \\
\midrule
\textbf{ciliated cells} 纤毛细胞 & move \textbf{mucus} 黏液 along the \textbf{trachea} 气管 and \textbf{bronchi} 支气管 & tiny hairs (cilia) on top sweep the mucus along \\
\textbf{root hair cells} 根毛细胞 & \textbf{absorption} 吸收 of water and minerals from the soil & a long, thin "hair" gives a large surface \\
\textbf{palisade mesophyll cells} 栅栏叶肉细胞 & photosynthesis & packed with chloroplasts, near the top of the leaf \\
\textbf{neurones} 神经元 & carry \textbf{electrical impulses} 电脉冲 & very long, to reach far across the body \\
\textbf{red blood cells} 红细胞 & transport \textbf{oxygen} 氧气 & no nucleus and a dish shape, to carry more oxygen \\
\textbf{sperm} 精子 and \textbf{egg cells} 卵细胞 & reproduction & the sperm can swim; the egg cell stores food \\
\bottomrule
\end{tabular}}
\end{center}

Sperm and egg cells are the sex cells, called \textbf{gametes} 配子.

\section*{Levels of organisation}

An \textbf{organism} 生物体 made of many cells is built up in levels, from small to large:

\begin{itemize}
\item \textbf{cell} — the smallest unit of life (for example a red blood cell).

\item \textbf{tissue} 组织 — a group of similar cells that work together (for example \textbf{muscle} 肌肉).

\item \textbf{organ} 器官 — several different tissues that work together to do a job (for example the heart, or a leaf).

\item \textbf{organ system} 器官系统 — several organs that work together (for example the \textbf{digestive system} 消化系统, which breaks down food).

\item \textbf{organism} — all the organ systems together make one whole living thing.

\end{itemize}

\section*{The size of cells: magnification}

Cells are tiny, so you look at them under a \textbf{microscope} 显微镜, which makes them appear much larger. How many times larger the image is, is called the \textbf{magnification} 放大倍数.

\[
\text{magnification} = \frac{\text{image size}}{\text{actual size}}
\]

\begin{itemize}
\item \textbf{image size} = the size in the drawing or photo.

\item \textbf{actual size} = the real size of the \textbf{specimen} 标本.

\end{itemize}

Magnification has \textbf{no unit} — it is just a number, for example ×250. You can rearrange the formula:

\[
\text{actual size} = \frac{\text{image size}}{\text{magnification}}, \qquad \text{image size} = \text{actual size} \times \text{magnification}
\]

\textbf{Worked example.} A structure is shown at a magnification of 250. Its image size is 5.00 mm. So the actual size = 5.00 ÷ 250 = 0.02 mm. Always put both sizes in the \textbf{same unit} before you divide.

\subsection*{Changing units (Supplement)}

Cells are often measured in \textbf{micrometres} 微米 (μm), which are smaller than millimetres (mm):

\[
1\ \text{mm} = 1000\ \mu\text{m}, \qquad 1\ \mu\text{m} = 0.001\ \text{mm}
\]

So 0.02 mm = 20 μm. To change mm into μm, multiply by 1000. To change μm into mm, divide by 1000.

\section*{Exam tips}

\begin{itemize}
\item A plant cell has a cell wall, a large vacuole and chloroplasts; an animal cell has none of these. Both have a nucleus, cytoplasm, cell membrane, mitochondria and ribosomes.

\item A cell with \textbf{no nucleus} but with plasmids and a loop of DNA is a bacterium.

\item For magnification, cover the quantity you want: magnification = image ÷ actual; actual = image ÷ magnification.

\item Always convert to the same unit before you calculate. Remember 1 mm = 1000 μm.

\item For a specialised cell, link its job to the one feature that suits it (for example a red blood cell has no nucleus, so it can hold more oxygen).

\end{itemize}


\end{document}
